\chapter{Định nghĩa hình thức các độ đo đánh giá}
\label{app:metrics}

Phụ lục này trình bày định nghĩa hình thức của các độ đo được dùng ở
Mục~\ref{sec:exp-metrics}. Các định nghĩa dạng văn xuôi đã được nêu trong phần thân;
phần dưới đây bổ sung công thức tường minh để tiện tra cứu.

Với hỏi đáp đọc hiểu, F1 ở cấp token là trung bình điều hòa của độ chính xác $P$ và độ
bao phủ $R$ giữa câu trả lời dự đoán và đáp án chuẩn, như Công thức~\eqref{eq:f1}:
\begin{equation}
  \label{eq:f1}
  \mathrm{F1} = \frac{2PR}{P + R}.
\end{equation}

Hệ số tương quan Matthews cho phân loại nhị phân, tính từ các đại lượng dương thật ($TP$),
âm thật ($TN$), dương giả ($FP$) và âm giả ($FN$), như Công thức~\eqref{eq:mcc}:
\begin{equation}
  \label{eq:mcc}
  \mathrm{MCC} = \frac{TP\cdot TN - FP\cdot FN}
  {\sqrt{(TP+FP)(TP+FN)(TN+FP)(TN+FN)}}.
\end{equation}

Với truy hồi và xếp hạng lại, gọi $\mathrm{rel}_i \in \{0,1\}$ là mức liên quan của kết quả ở
hạng $i$ và $\mathrm{rank}_q$ là thứ hạng của ngữ cảnh đúng đầu tiên cho truy vấn $q$, hai độ
đo được xác định như Công thức~\eqref{eq:ndcg-mrr}:
\begin{equation}
  \label{eq:ndcg-mrr}
  \mathrm{nDCG@}k = \frac{1}{Z}\sum_{i=1}^{k}\frac{2^{\mathrm{rel}_i}-1}{\log_2(i+1)},
  \qquad
  \mathrm{MRR} = \frac{1}{|Q|}\sum_{q\in Q}\frac{1}{\mathrm{rank}_q},
\end{equation}
trong đó $Z$ là hệ số chuẩn hóa (DCG của thứ hạng lý tưởng) đưa nDCG về đoạn $[0,1]$, và $Q$
là tập truy vấn đánh giá.
